\section{Conclusions and outlook}
\label{sec:conclusion}

In this study, we theoretically explored the role of local environments (or the short-range chemical disorder), specifically atomic configurations and compositions, on the intrinsic Gilbert damping. Our investigation is focused on various bcc FeCo alloys, namely Fe$_{100-x}$Co$_{x}$, with $x\in[0\%,60\%]$. With that perspective, a hybrid explicit/effective medium model (EC-VCA) was introduced. 
%The results underscore that local atomic arrangements markedly affect the damping. 
From the analysis of the complete set of configurations in the Fe$_{87}$Co$_{13}$ concentration, it was demonstrated that different spatial distributions of Fe/Co within the embedded cluster region can markedly affect the damping. Unlike other quantities such as the Heisenberg exchange interactions, damping is anisotropic with respect to the spin quantization axis, which further contributes (albeit slightly in FeCo) to the effects of chemical disorder. The computed average damping aligns well with both experiments and the pure VCA calculation. When low-temperature explicit atomistic spin dynamics simulations are performed, the influence of short-range disorder on local dynamics is observed to alter the overall relaxation rate in both nonlocal and effective damping formulations of the generalized Landau-Lifshitz-Gilbert (LLG) equation, even when such a small (15-atom) embedded region is considered.

%We first investigated the complete set of configurations in Fe$_{87}$Co$_{13}$, using clusters embedded in a VCA matrix, and found that the damping variation with respect to the spin quantization axis is around 1\%. Additionally, its average is higher in the cluster based calculations compared to pure VCA results of the damping, which is due to the enhanced electron lifetimes.

Transitioning to configurations with higher Co content across the FeCo alloy series, for which inevitably only representative sets can be analyzed, the effective damping is found to be notably influenced by the atomic configuration at the nearest neighborhood. In the most extreme cases, a variation of about one order of magnitude is found. A direct correlation between effective damping and average the Co interatomic distance ($d$), as well as the quantity of Co atoms occupying the first neighborhood ($n$) were found. This local environment-dependent damping behavior was explained in light of the simplified Kambersky’s formula (Eq. \ref{eq:kambersky}), demonstrating a consistent correlation between damping and the local density of states at the Fermi level. In a global perspective (i.e., performing a configuration average), those differences in damping are masked by statistical averages, reconciling the local picture to what is typically observed in FMR experiments for FeCo.
%Lastly, we compare the average damping of selected clusters within each composition with the experiment as well as alloy theory. 

Considering damping as a nonlocal (and inhomogeneous) quantity reveals several interesting characteristics, such as a pronounced dependence on the local environment, which are overlooked by effective medium theories. Such local environment-dependent damping may have implications, in real materials, on phenomena like the damping anisotropy observed in Fe$_{50}$Co$_{50}$, which was originally attributed to short-range configurational effects  \cite{PhysRevLett.122.117203}. Moreover, our results naturally encourage further exploration of the potential for local engeneering of damping, and the investigation into how defects affect both damping and the magnetization dynamics at local and global scales. 
%Our findings suggest that the damping parameter in the Landau-Lifshitz-Gilbert (LLG) equation should not be viewed as a homogeneous attribute in magnetic alloys. Instead, it exhibits a dependence on the local atomic environment, similar to the dependence of local environment on the magnetic anisotropy [47] or magnetic moment. Such local environment-dependent damping has implications for magnetization dynamics and magnon properties in materials, with an anticipated pronounced effect in amorphous materials due to their highly localized environments. This study not only sheds light on the origins of damping inhomogeneity in magnetic materials but also opens avenues for manipulating damping through atomic engineering, presenting significant implications for material design and applications.

